\documentclass[12pt]{article}
\usepackage{notes-common}
\usetikzlibrary{calc}
\geometry{margin=0.35in}
\AtBeginDocument{\fontsize{14}{17}\selectfont}

% Complete / filled-in reference version of circuits-intro-notes.tex.
% Mon Oct 5 -- Ohm's law + Kirchhoff's laws (Knight 27.5, 28.1-28.2).
% Built from Seth's rough notes (uploaded).
\newcommand{\ANS}[1]{\textcolor{accent}{#1}}

% ---- same battery/resistor schematic pieces as the skeleton ----
\newcommand{\battv}[1]{% vertical battery symbol centered at (#1)
  \draw[line width=1.1pt] ($(#1)+(-0.35,0.12)$) -- ($(#1)+(0.35,0.12)$);
  \draw[line width=2.4pt] ($(#1)+(-0.2,-0.12)$) -- ($(#1)+(0.2,-0.12)$);
}
\newcommand{\reszig}[1]{% horizontal resistor zigzag centered at (#1)
  \fill[white] ($(#1)+(-0.65,-0.32)$) rectangle ($(#1)+(0.65,0.32)$);
  \draw ($(#1)+(-0.6,0)$)
    -- ($(#1)+(-0.42,0.22)$) -- ($(#1)+(-0.24,-0.22)$)
    -- ($(#1)+(-0.06,0.22)$) -- ($(#1)+(0.12,-0.22)$)
    -- ($(#1)+(0.3,0.22)$) -- ($(#1)+(0.48,-0.22)$)
    -- ($(#1)+(0.6,0)$);
}

\begin{document}

\noteshead{Battery \& Resistor Circuits \textcolor{accent}{(completed)}}{}

\goal{
\begin{itemize}[itemsep=5pt,parsep=2pt,topsep=2pt]
\item Pin down what current and $\Delta V$ actually mean, and the four
rules that follow just from charge and energy conservation.
\item Turn Ohm's law into a tool: given a resistor and a $\Delta V$ across
it, get $I$.
\item Work a single resistor both with an ideal battery and with a real
one (internal resistance $r$), to see what $r$ costs you.
\item Derive $R_S=R_1+R_2$ for series and $R_P=(1/R_1+1/R_2)^{-1}$ for
parallel from the rules in \S1 --- not as formulas to memorize, but as
consequences of ``same current'' and ``same $\Delta V$.''
\end{itemize}}

%============================================================
\nsec{Definitions}

\textbf{Current}: \ANS{the amount of charge passing a point (through an
element) per second.} \quad $I\equiv\Delta q/\Delta t$.

\hspace{1cm} units: \ANS{1 A $= 1$ C/s}

\textbf{Change in voltage, $\Delta V$}: \ANS{the potential energy per
charge gained or lost by a charge that crosses an element.}

\vspace{2mm}
\textbf{Series}: elements connected \ANS{back-to-back --- one shared wire
between them, so the same charge has to pass through each in turn.}

\textbf{Parallel}: elements connected \ANS{side-by-side, sharing \emph{both}
endpoints (the same two nodes).}

%============================================================
\nsec{Four rules}

\begin{itemize}[itemsep=7pt,parsep=3pt,topsep=3pt]
\item Elements in series have \ANS{the same current passing through them.}
\item Elements in parallel have \ANS{the same voltage change ($\Delta V$)
across them.}
\item Current is \ANS{conserved} at a junction.
\item Around any closed loop, the $\Delta V$'s \ANS{sum to zero.}
\end{itemize}

\checkpoint{The last two rules are each a conservation law in disguise. Which
one is charge conservation, and which one is energy conservation?}

{\small\color{accent} The junction rule (current conserved) \emph{is} charge
conservation --- charge can't pile up or vanish at a junction, so whatever
flows in must flow out. The loop rule ($\Delta V$'s sum to zero) is energy
conservation --- walk all the way around back to where you started, and the
net PE per charge you gained and lost along the way has to cancel.}

%============================================================
\nsec{Ohm's law defines a resistor}

\[
\Delta V = \ANS{IR}
\]
\vspace{0.3cm}

Given a resistor of resistance $R$ and a voltage difference $\Delta V$
across it, the current through it is

\[
I = \ANS{\dfrac{\Delta V}{R}}
\]

%============================================================
\nsec{One resistor and a battery}

\begin{center}
\begin{tikzpicture}[x=1cm,y=1cm]
  \draw (0,0) -- (0,2) -- (3,2) -- (3,0) -- cycle;
  \battv{0,1}
  \node[scale=0.8, anchor=east] at (-0.15,1) {$\varepsilon,\,r$};
  \reszig{1.5,2}
  \node[scale=0.8, anchor=south] at (1.5,2.45) {$R$};
\end{tikzpicture}
\end{center}

\textbf{(a) Ideal battery}: $\varepsilon=12$ V, $r=0$, $R=4\ \Omega$.

\[
I = \dfrac{\varepsilon}{R} = \dfrac{12}{4} = \ANS{3.0\text{ A}}
\qquad
V_B = \varepsilon - Ir = 12-0 = \ANS{12\text{ V}}
\]

\textbf{(b) Real battery}: same $\varepsilon$ and $R$, but now $r=1\ \Omega$.

\[
I = \dfrac{\varepsilon}{R+r} = \dfrac{12}{5} = \ANS{2.4\text{ A}}
\qquad
V_B = \varepsilon - Ir = 12-(2.4)(1) = \ANS{9.6\text{ V}}
\]

\checkpoint{$V_B$ dropped in part (b). Where did that energy per charge go?}

{\small\color{accent} Into the battery's own internal resistance $r$ ---
the missing $12-9.6=2.4$ V is exactly $Ir=(2.4\text{ A})(1\ \Omega)=2.4$ V.
That energy per charge is dissipated \emph{inside} the battery itself, not
delivered to $R$.}

%============================================================
\nsec{Two resistors in series}

\begin{center}
\begin{tikzpicture}[x=1cm,y=1cm]
  \draw (0,0) -- (0,2) -- (6,2) -- (6,0) -- cycle;
  \battv{0,1}
  \node[scale=0.8, anchor=east] at (-0.15,1) {$\varepsilon$};
  \reszig{2,2}
  \node[scale=0.8, anchor=south] at (2,2.45) {$R_1$};
  \reszig{4,2}
  \node[scale=0.8, anchor=south] at (4,2.45) {$R_2$};
\end{tikzpicture}
\end{center}

\textbf{Derive $R_S$} (symbols only --- use the series current rule and Ohm's
law on each resistor, then the loop rule):

{\small
Series rule: the same current $I$ flows through both resistors. Ohm's law
on each: $\Delta V_1 = IR_1$, $\Delta V_2 = IR_2$. Loop rule:
\[
\varepsilon = \Delta V_1+\Delta V_2 = IR_1+IR_2 = I(R_1+R_2).
\]
But by definition of the single equivalent resistor, $\varepsilon = IR_S$.
Comparing the two:}

\[
\boxed{R_S = R_1+R_2}
\]

\textbf{Now with numbers}: $\varepsilon=12$ V, $R_1=4\ \Omega$, $R_2=8\ \Omega$
(ideal battery).

\[
R_S = 4+8 = \ANS{12\ \Omega}
\qquad
I = \dfrac{\varepsilon}{R_S} = \dfrac{12}{12} = \ANS{1.0\text{ A}}
\]
\[
\Delta V_1 = IR_1 = \ANS{4.0\text{ V}}
\qquad
\Delta V_2 = IR_2 = \ANS{8.0\text{ V}}
\]

\checkpoint{Check the loop rule with your two $\Delta V$'s above: do they
add up to $\varepsilon$?}

{\small\color{accent} Yes: $4.0+8.0=12$ V $=\varepsilon$. That's not a
coincidence --- it's exactly the loop rule, built into how $R_S$ was
derived.}

%============================================================
\nsec{Two resistors in parallel}

\begin{center}
\begin{tikzpicture}[x=1cm,y=1cm]
  \draw (0,0) -- (0,2) -- (1.3,2) -- (1.3,2.9) -- (4.7,2.9) -- (4.7,2) -- (6,2) -- (6,0) -- cycle;
  \draw (1.3,2) -- (1.3,1.1) -- (4.7,1.1) -- (4.7,2);
  \battv{0,1}
  \node[scale=0.8, anchor=east] at (-0.15,1) {$\varepsilon$};
  \reszig{3,2.9}
  \node[scale=0.8, anchor=south] at (3,3.35) {$R_1$};
  \reszig{3,1.1}
  \node[scale=0.8, anchor=north] at (3,0.65) {$R_2$};
\end{tikzpicture}
\end{center}

\textbf{Derive $R_P$} (symbols only --- use the parallel voltage rule and
Ohm's law on each resistor, then the junction rule at the top node):

{\small
Parallel rule: both resistors see the same $\Delta V$ (here, $\Delta
V=\varepsilon$, since each branch connects straight across the ideal
battery). Ohm's law on each: $I_1=\Delta V/R_1$, $I_2=\Delta V/R_2$.
Junction rule at the top node:
\[
I = I_1+I_2 = \dfrac{\Delta V}{R_1}+\dfrac{\Delta V}{R_2}
= \Delta V\left(\dfrac1{R_1}+\dfrac1{R_2}\right).
\]
But by definition of the single equivalent resistor, $I=\Delta V/R_P$.
Comparing the two:}

\[
\boxed{R_P = \left(\dfrac1{R_1}+\dfrac1{R_2}\right)^{-1}}
\]

\textbf{Now with numbers}: $\varepsilon=12$ V, $R_1=4\ \Omega$,
$R_2=12\ \Omega$ (ideal battery).

\[
R_P = \left(\dfrac14+\dfrac1{12}\right)^{-1} = \left(\dfrac{3}{12}+\dfrac1{12}\right)^{-1}
= \left(\dfrac{4}{12}\right)^{-1} = \ANS{3.0\ \Omega}
\]
\[
I_1 = \dfrac{\varepsilon}{R_1} = \dfrac{12}{4} = \ANS{3.0\text{ A}}
\qquad
I_2 = \dfrac{\varepsilon}{R_2} = \dfrac{12}{12} = \ANS{1.0\text{ A}}
\]
\[
I_{\text{total}} = I_1+I_2 = \ANS{4.0\text{ A}}
\quad\big(\text{check: } \varepsilon/R_P = 12/3 = 4.0\text{ A}\ \checkmark\big)
\]

\checkpoint{Which resistor carries more current, and why does that make
sense given which one is smaller?}

{\small\color{accent} $R_1$ ($4\ \Omega$) carries more current ($3.0$ A vs.\
$1.0$ A). Both resistors feel the \emph{same} $\Delta V$, and $I=\Delta
V/R$ --- so the smaller resistance lets more current through for that same
push. Smaller $R$ means less obstruction, not less current.}

\end{document}
